\section{A mixed finite cover lifts a cross-cup arithmetic character}
\label{sec:cross-cup-five-twenty-nine-lift}

Let \(G_0\) be the positive peripheral automorphism group
containing the profinite Grothendieck--Teichm\"uller group
\(G\). Its action on \(F=\widehat F(x,y)\) has the form
\begin{equation}
 \phi_g(x)=x^{\lambda_g},\qquad
 \phi_g(y)=f_g^{-1}y^{\lambda_g}f_g,\qquad
 \phi_g(z)\sim z^{\lambda_g},\qquad z=(xy)^{-1}.
                                                               \label{eq:cross-cup-5-29-1}
\end{equation}
Here conjugacy in the last formula is in \(F\), and the
same multiplier occurs on all three peripheral generators.
Assume \(f_g\in F\,{}^{\prime}\), as in the standard GT action. Put \(M_0=\ker(\lambda:G_0\to\widehat{\mathbf Z}^{\times})\).

\begin{theorem}\label{thm:five-twenty-nine-cross-cup-lift}
There is a continuous \(G_0\)-invariant homomorphism
\begin{equation}
                    \xi_{5,29}:M_0\longrightarrow\mathbf F_2
                                                               \label{eq:cross-cup-5-29-2}
\end{equation}
whose restriction to \(G_K\), for
\(K=\mathbf Q(\mu_\infty,2^{1/\infty})\), is precisely the
central quadratic character \(\kappa_{5,29}\) of
\[
 \mathbf Q\!\left(\sqrt5,\sqrt{29},\sqrt{7+2\sqrt5}\right).
\]
In particular its restriction to
\(N=\ker(\lambda,\rho_2)\subset G\) belongs to
\(X_2=H^1(N,\mathbf F_2)^G\), and its arithmetic restriction
is outside the span of the rational quartic Kummer
characters and \(\eta_3\).

Equivalently, the corrected cross-cup class
\begin{equation}
                         e_{c_5,c_{29}}=c_5\cup c_{29}
                                                               \label{eq:cross-cup-5-29-3}
\end{equation}
has zero inflation to \(H^2(G,\mathbf F_2)\).
The construction uses the homology modulo twenty-nine of
a cyclic degree-five cover, and applies already to \(G_0\).
\end{theorem}

\subsection{The cyclic cover and its complete peripheral pairing}

Let
\begin{equation}
 U_5=\ker\bigl(F\longrightarrow C_5,\quad x\mapsto1,\ y\mapsto1\bigr).
                                                               \label{eq:cross-cup-5-29-4}
\end{equation}
Every \(\phi_g\) preserves \(U_5\), because it multiplies
the indicated abelian quotient by
\(\lambda_g\bmod5\). Put
\begin{equation}
 V=H_1(U_5,\mathbf F_{29})/
             \langle\text{peripheral homology classes}\rangle .
                                                               \label{eq:cross-cup-5-29-5}
\end{equation}
This definition uses the subgroup action in \eqref{eq:cross-cup-5-29-1} directly;
it does not assume equivariance of any additional
covering or filling operation.

\begin{lemma}\label{lem:five-cover-similitude}
The vector space \(V\) has dimension four and carries a
nondegenerate alternating form.
The induced representation
\begin{equation}
                    \rho:G_0\longrightarrow
                                  \operatorname{GSp}(V)
                                                               \label{eq:cross-cup-5-29-6}
\end{equation}
has similitude multiplier \(\lambda_g\bmod29\).
If \(D\) is the deck transformation induced by conjugation
by \(x\), then
\begin{equation}
                  \rho(g)D\rho(g)^{-1}=D^{\lambda_g\bmod5}.
                                                               \label{eq:cross-cup-5-29-7}
\end{equation}
The characteristic polynomial of \(D\) on \(V\) is
\(1+X+X^2+X^3+X^4\).
\end{lemma}

\begin{proof}
Start with the discrete oriented pair of pants.
The degree-five cover specified by \eqref{eq:cross-cup-5-29-4} has Euler
characteristic \(-5\).
Each of \(x,y,z\) maps to a generator of \(C_5\)
(their images are \(1,1,-2\)), so the cover has exactly
three punctures.
Its compactification has genus two.
Consequently \(U_5\) is the profinite completion of a
free group of rank six with a surface presentation
\begin{equation}
 [a_1,b_1][a_2,b_2]c_1c_2c_3=1.                           \label{eq:cross-cup-5-29-8}
\end{equation}
The three peripheral classes have the sole homology
relation \(c_1+c_2+c_3=0\).
Their span has dimension two, and \eqref{eq:cross-cup-5-29-5} has dimension four.

There is a single cusp above each original cusp.
Thus \eqref{eq:cross-cup-5-29-1}, restricted to \(U_5\), sends each \(c_i\) to
a \(U_5\)-conjugate of \(c_i^{\lambda_g}\).
To check the conjugator assertion, a conjugator in \(F\)
may be multiplied by a power of the relevant original
peripheral generator to move it into \(U_5\), since that
peripheral generator maps onto the deck group; this
multiplication does not change its conjugation on \(c_i\).

We prove the multiplier assertion algebraically.
Take the class-two, exponent-twenty-nine quotient of
\(U_5\).
Its commutator space is the exterior square of its
first homology.
Project that exterior square onto \(\bigwedge^2 V\).
Conjugating a peripheral word changes its central
coordinate by a wedge with its peripheral homology
class, which vanishes under this projection.
Taking a \(\lambda_g\)-th power multiplies its projected
central coordinate by \(\lambda_g\bmod29\).
Relation \eqref{eq:cross-cup-5-29-8} therefore implies
\begin{equation}
 \bigwedge\nolimits^2\rho(g)
       (a_1\wedge b_1+a_2\wedge b_2)
   =\lambda_g
       (a_1\wedge b_1+a_2\wedge b_2)
       \quad\text{in }\bigwedge\nolimits^2V.               \label{eq:cross-cup-5-29-9}
\end{equation}
Here the displayed \(a_i,b_i\) mean their homology classes.
For an explicit way to read this calculation, eliminate
\(c_3\) from \eqref{eq:cross-cup-5-29-8}, use
\(a_1,b_1,a_2,b_2,c_1,c_2\) as free generators, and take
class-two logarithms, which are valid because twenty-nine
is odd. The projected central coordinate of \(c_3\) is
the negative of the bivector in \eqref{eq:cross-cup-5-29-9}, while those of
\(c_1,c_2\) are zero. The peripheral conjugation and
power rule gives exactly \eqref{eq:cross-cup-5-29-9}.

That bivector is nondegenerate. Its inverse alternating
form consequently has the same similitude multiplier.
This proves \eqref{eq:cross-cup-5-29-6}, including the full multiplier statement.

Conjugation by \(x\) induces the deck action, and
\(\phi_g(x)=x^{\lambda_g}\) proves \eqref{eq:cross-cup-5-29-7}.
Finally the cellular chain complex of the five-sheeted
cover of the two-generator graph has
\[
 C_1=\mathbf F_{29}[C_5]^2,\qquad C_0=\mathbf F_{29}[C_5].
\]
At every nontrivial deck character \(\zeta^j\), its
boundary map is the nonzero row
\((\zeta^j-1,\zeta^j-1)\); its kernel is one-dimensional.
The trivial character has multiplicity two in homology,
and is exactly the peripheral subspace.
Thus each nontrivial fifth-root character occurs once
in \(V\), proving the characteristic polynomial.
\end{proof}

\subsection{The finite symplectic normalizer}

Write \(q=29\), \(k=\mathbf F_{q^2}\), and
\(\bar z=z^q\). Choose \(\zeta\in k\) of order five.
Since \(q\equiv-1\pmod5\), the four deck characters
split into the two Frobenius pairs
\(\{\zeta,\zeta^{-1}\}\) and
\(\{\zeta^2,\zeta^{-2}\}\).
We may identify \(V\) with \(k\oplus k\), with
\begin{equation}
                         D(z_1,z_2)=(\zeta z_1,\zeta^2z_2).
                                                               \label{eq:cross-cup-5-29-10}
\end{equation}
The two summands are orthogonal and nondegenerate.
On each, a deck-invariant alternating form is
\(\operatorname{Tr}_{k/\mathbf F_q}(d_i z\bar w)\),
with \(d_i\ne0\) and \(\bar d_i=-d_i\).
Rescaling the coordinate by an element of prescribed
norm allows taking \(d_1=d_2=d\).
The full form is then
\begin{equation}
 \langle z,w\rangle
 =\operatorname{Tr}_{k/\mathbf F_q}
       \bigl(d(z_1\bar w_1+z_2\bar w_2)\bigr),
 \qquad \bar d=-d.                                        \label{eq:cross-cup-5-29-11}
\end{equation}

Let
\[
                    P(z_1,z_2)=(z_2,\bar z_1).
\]
Then \(P^4=1\) and \(PDP^{-1}=D^2\).
Every linear map normalizing \(\langle D\rangle\)
is uniquely of the form \((v_1,v_2)P^j\), with
\(v_i\in k^\times\) and \(j\in\mathbf Z/4\).
Let \(E\) be the subgroup of such maps which are
similitudes of \eqref{eq:cross-cup-5-29-11}, and write
\[
 \pi:E\longrightarrow C_4\times\mathbf F_q^\times
\]
for deck exponent \(2^j\) and similitude multiplier.
Its kernel is
\begin{equation}
 T=\{(v_1,v_2):N(v_1)=N(v_2)=1\}
                   \simeq C_{30}\times C_{30}.             \label{eq:cross-cup-5-29-12}
\end{equation}
Define
\begin{equation}
                 (-1)^{\chi(v_1,v_2)}=(v_1v_2)^{15}.
                                                               \label{eq:cross-cup-5-29-13}
\end{equation}
This is a nonzero homomorphism \(T\to\mathbf F_2\).

\begin{lemma}\label{lem:finite-cross-cup-extension}
The subgroup \(\ker\chi\) is normal in \(E\), and
\begin{equation}
 1\longrightarrow C_2\longrightarrow E/\ker\chi
       \longrightarrow C_4\times\mathbf F_{29}^\times
       \longrightarrow1                                 \label{eq:cross-cup-5-29-14}
\end{equation}
is a central extension with class \(x\cup y\).
Here \(x\) is the nonzero quadratic character of \(C_4\),
and \(y\) is the quadratic character of
\(\mathbf F_{29}^\times\).
\end{lemma}

\begin{proof}
Conjugation by the scalar subgroup is trivial on \(T\).
Conjugation by \(P\) sends \((v_1,v_2)\) to
\((v_2,\bar v_1)=(v_2,v_1^{-1})\) on \(T\).
This leaves \eqref{eq:cross-cup-5-29-13} unchanged, because \(v_1^{15}\)
has order at most two.
Thus \(\ker\chi\) is normal and its quotient kernel is
central.

Choose \(\beta\in k^\times\) with \(N\beta=-1\), and set
\[
                          J=(1,\beta)P.
\]
From \eqref{eq:cross-cup-5-29-11}, \(J\) is a similitude of multiplier one:
\(P\) interchanges the two summands and reverses the
sign of the second contribution, which multiplication
by \(\beta\) reverses again.
More explicitly, a map \((v_1,v_2)P\) has multiplier
\(t\) precisely when \(Nv_1=t\) and \(Nv_2=-t\).
Thus \(\pi(J)=(1,1)\).

Choose \(v\in k^\times\) with norm a generator \(t\)
of \(\mathbf F_{29}^\times\), and put
\[
                             L=(v,v).
\]
Then \(\pi(L)=(0,t)\).
These two images generate the entire base of \eqref{eq:cross-cup-5-29-14};
in particular \(\pi\) is surjective.

There are three independent degree-two classes on
\(C_4\times C_{28}\): the two cyclic extension classes
and the mixed cup product \(xy\).
Their coefficients are determined respectively by
the fourth power of a lift of the first generator,
the twenty-eighth power of a lift of the second,
and their commutator.
We compute all three in the central quotient.

First,
\[
                       J^4=(-1,-1).
\]
This need not be the identity
in \(E\).
Nevertheless
\[
                       \chi(J^4)=0,
\]
because the two norm-one quadratic signs cancel.
Next,
\[
                       L^{28}=(v^{28},v^{28})\in T,
 \qquad                 \chi(L^{28})=0.
\]
Finally,
\[
                       [J,L]=(1,\bar v/v)\in T,
\]
and
\begin{equation}
 \begin{split}
 (-1)^{\chi([J,L])}
   &=(\bar v/v)^{15}
     =v^{420}
     =(Nv)^{14}
     =t^{14}
     =-1 .
 \end{split}                                               \label{eq:cross-cup-5-29-15}
\end{equation}
Thus both cyclic coefficients vanish and the mixed
coefficient is one. This proves the claimed exact
extension class.
\end{proof}

\subsection{The full GT lift and its explicit beta formula}

By Lemma~\ref{lem:five-cover-similitude}, the representation
\(\rho\) takes values in \(E\), and
\begin{equation}
                   \pi\rho(g)=(\lambda_g\bmod5,\lambda_g\bmod29).
                                                               \label{eq:cross-cup-5-29-16}
\end{equation}
The identification of \((\mathbf Z/5)^\times\) with
\(C_4\) uses its generator \(2\).
Under \eqref{eq:cross-cup-5-29-16}, the characters \(x,y\) in
Lemma~\ref{lem:finite-cross-cup-extension} pull back
to \(c_5,c_{29}\), respectively.
Thus \(\rho\), followed by \(E\to E/\ker\chi\), is a
full-\(G_0\) lift of the cyclotomic finite quotient
through the central extension of class \(c_5c_{29}\).
This proves its vanishing in \(H^2(G_0,\mathbf F_2)\),
and hence in \(H^2(G,\mathbf F_2)\).

For \(g\in M_0\), both components of \eqref{eq:cross-cup-5-29-16} are trivial,
so \(\rho(g)\in T\).
Define
\[
                          \xi_{5,29}(g)=\chi(\rho(g)).
\]
Normality and centrality in the preceding lemma prove
that this is a continuous \(G_0\)-invariant homomorphism.

The construction also gives a direct beta expression.
Write \(B_g(X,Y)\) for the normalized Alexander coefficient
in the usual convention: when \(\lambda_g=1\), it is the
scalar action on the nontrivial character line of
\(F'/F''\).
Then
\begin{equation}
 (-1)^{\xi_{5,29}(g)}
  =\left(
      B_g(\zeta,\zeta)
      B_g(\zeta^2,\zeta^2)
    \right)^{15}\qquad(g\in M_0),                          \label{eq:cross-cup-5-29-17}
\end{equation}
with coefficients reduced modulo twenty-nine.
The commutator \([x,y]\) projects nontrivially to each
of these character lines: in the graph chain complex
its value at \((\zeta^j,\zeta^j)\) is a nonzero multiple
of the difference of the two edge classes.
Thus the scalar in the homology construction is indeed
the indicated beta value, without a change-of-basis
factor.
The norm-one property follows from \eqref{eq:cross-cup-5-29-11} and multiplier
one.
Replacing \(\zeta\) by another primitive fifth root
permutes the two factors with possible inversions,
which leave their quadratic sign unchanged.
For the Alexander convention, compare
\cite[\S2.3, p.~101]{Furusho2003}; the finite homology
construction above is also an independent definition
of \eqref{eq:cross-cup-5-29-2}.

Apply naturality of transgression first to
\(1\to M_0\to G_0\to\lambda(G_0)\to1\), and then to
\(1\to N\to G\to A\to1\).
The central extension \eqref{eq:cross-cup-5-29-14} gives
\begin{equation}
                         d_G(\xi_{5,29}|_N)=c_5c_{29}
                              \quad\text{in }H^2(A,\mathbf F_2).
                                                               \label{eq:cross-cup-5-29-18}
\end{equation}
Here \eqref{eq:cross-cup-5-29-18} can equivalently be read by pulling back \eqref{eq:cross-cup-5-29-14}
along \(A\to U\to C_4\times\mathbf F_{29}^\times\);
the map from \(G\) to that pullback has central restriction
\(\xi_{5,29}|_N\).

On the arithmetic subgroup, the independently computed
dihedral character \(\kappa_{5,29}\) has the same
transgression \(c_5c_{29}\).
Its Hilbert corrections vanish because
\[
                         N_{\mathbf Q(\sqrt5)/\mathbf Q}
                                  (7+2\sqrt5)=29.
\]
The arithmetic transgression is injective, since \(K\)
contains \(\mathbf Q^{\mathrm{ab}}\).
Therefore
\[
                         \xi_{5,29}|_{G_K}=\kappa_{5,29}.
\]
This proves Theorem~\ref{thm:five-twenty-nine-cross-cup-lift}.

\subsection{A further strict normalization}

Retain the established factorization \(G=R_DH\), the
residual subgroup \(V_D\), and the established
fact that \(R_D/V_D\) is pro-odd.
Define
\begin{equation}
 R_D^\xi=R_D\cap\ker\xi_{5,29},\qquad
 V_D^\xi=V_D\cap\ker\xi_{5,29}.                            \label{eq:cross-cup-5-29-19}
\end{equation}
Restrictions of \(\xi_{5,29}\) to these groups are
understood through their inclusions in \(M_0\).

\begin{corollary}\label{cor:cross-cup-strict-normalization}
The subgroups in \eqref{eq:cross-cup-5-29-19} are closed and \(G\)-normal, and
\begin{equation}
 [R_D:R_D^\xi]=[V_D:V_D^\xi]=2,\qquad
                         G=R_D^\xi H.                     \label{eq:cross-cup-5-29-20}
\end{equation}
The maximal pro-two relative quotient is unchanged
on passage to \(V_D^\xi\):
\begin{equation}
 B_{\max2}\simeq
 \left(
 V_D^\xi/
 \overline{[G,V_D^\xi](V_D^\xi\cap H)}
 \right)_{\max2}.                                         \label{eq:cross-cup-5-29-21}
\end{equation}
If the established kernel of restriction from \(X_2\)
to \(H^1(V_D,\mathbf F_2)\) is the span of the rational
Kummer characters and \(\eta_3\), then the kernel of
restriction to \(H^1(V_D^\xi,\mathbf F_2)\) is exactly
that span plus the independent character
\(\xi_{5,29}|_N\).
No factorization \(G=V_D^\xi H\) is asserted.
\end{corollary}

\begin{proof}
The arithmetic inertia commutator of
Proposition~\ref{prop:cross-cup-residual-inertia} belongs
to \(V_D\cap H\) and has \(\xi_{5,29}\)-value one.
This proves both indices in \eqref{eq:cross-cup-5-29-20}.
For any \(r\in R_D\), multiplying by that arithmetic
element when necessary gives an element of \(R_D^\xi\).
Together with \(G=R_DH\), this proves \(G=R_D^\xi H\).

The structural relative quotient may therefore be
computed using \(R_D^\xi\).
The quotient \(R_D^\xi/V_D^\xi\) embeds into the pro-odd
group \(R_D/V_D\), so the established central coprime
extension argument identifies their maximal pro-two
relative quotients. This proves \eqref{eq:cross-cup-5-29-21}, using a global
factorization at \(R_D^\xi\), not an unproved internal
factorization assertion at \(V_D^\xi\).

Finally, a character in \(X_2\) vanishing on \(V_D^\xi\)
restricts on \(V_D\) to a character of its quotient
\(C_2\).
Thus its restriction is either zero or
\(\xi_{5,29}|_{V_D}\).
Subtracting the latter when necessary gives precisely
the old restriction kernel.
The arithmetic commutator proves independence from
that old kernel.
\end{proof}
